\scanpage{164}
\unsure{construction}'' qui apparaissent

\diagram{\(\displaystyle
(264)\qquad\left\{\vcenter{\hbox{%
\begin{tikzcd}[ampersand replacement=\&,column sep=large,row sep=small]
\mathbb C^* \& \& \text{(Espace analytique ``global'')}\\
\mathbb P \arrow[u,bend left=55] \& \text{point} \& \text{(id)}\\
\mathbb D \& \text{germe de disque} \& \\
\mathbb D^* \& \text{germe de disque épointé} \&
\end{tikzcd}}}\right.\)}

dont deux ``globaux'', et deux ``locaux'', qui s'assemblent suivant le diagramme d'inclusions

\[
(265)
\]
\diagram{\begin{tikzcd}[ampersand replacement=\&,row sep=large,column sep=large]
\&\mathbb P \arrow[d,hook]\\
\mathbb D^* \arrow[r,hook] \arrow[d,hook]\&\mathbb D \arrow[d,dotted]\\
\mathbb C^* \arrow[r,dotted]\&\bigl[\mathbb C\bigr]
\end{tikzcd}}

La variété analytique globale $\mathbb C$, avec son diviseur $\unsure{\Theta|0}$ localisé en $\mathbb P=\{0\}$, se reconstitue à partir de ce diagramme comme somme amalgamée
\[
(266)\qquad \mathbb C\simeq \mathbb C^*\amalg_{\mathbb D^*}\mathbb D.
\]

Si on veut comprendre la situation ``de codim \unsure{un}'' dans les \unsure{représentants} typiques les plus élémentaires, on prendra
\margin{163}
